% ============================================================================
%  Subdocumento 10. Contenido. Se usa igual en el documento único y en el
%  PDF propio del subdocumento. Los formularios que le asigna el T-21 se
%  agregan solos al final (datos en formularios/ de esta carpeta).
% ============================================================================

\begin{resumenApertura}
\marcador{Qué resuelve este subdocumento, en cuatro o cinco líneas}
\begin{recibe}
  \item \marcador{Qué recibe el mandante}
\end{recibe}
\end{resumenApertura}

\section{Modelo de soporte}\label{sec:10-modelo-de-soporte}

\marcador{ITIL 4, niveles, canales, horarios y escalamiento}

\section{Indicadores y niveles de servicio}\label{sec:10-indicadores-y-niveles-de-servici}

\marcador{Coherentes con el Artículo 78}

\section{Dimensionamiento de la mesa de servicio}\label{sec:10-dimensionamiento-de-la-mesa-de-s}

\marcador{Teoría de colas, Erlang C u otro modelo declarado}

\section{Acuerdos operacionales y contratos de apoyo}\label{sec:10-acuerdos-operacionales-y-contrat}

\marcador{Coherencia con los compromisos externos}

\section{Libros de operación y conocimiento}\label{sec:10-libros-de-operacion-y-conocimien}

\marcador{Guías de resolución y automatización progresiva}

\section{Observabilidad de extremo a extremo}\label{sec:10-observabilidad-de-extremo-a-extr}

\marcador{Correlación de eventos y detección proactiva}
